\textbf{\textit{Soal. }}In an acute triangle $ABC$, let $M$ be the midpoint of $\overline{BC}$. Let $P$ be the foot of the perpendicular from $C$ to $AM$. Suppose that the circumcircle of triangle $ABP$ intersects line $BC$ at two distinct points $B$ and $Q$. Let $N$ be the midpoint of $\overline{AQ}$. Prove that $NB=NC$.

\begin{proof}[\textbf{Solusi 1.}] (Friday, 21 July 2023) Misalkan $D$ adalah proyeksi $A$ ke $BC$. Berarti kita punya $\dangle CPA = \dangle CDA = 90^\circ$ yang langsung mengakibatkan $CPDA$ siklis. 

\begin{center}
    \begin{asy}
        import olympiad;
        import geometry;
        size(250);
        pair A,B,C,M,P,Q,N,D;
        A=dir(75);
        B=dir(210);
        C=dir(-30);
        M=(B+C)/2;
        P=foot(C,A,M);
        circle c = circumcircle(triangle(A,B,P));
        pair BQ[] = intersectionpoints(c, line(B,C));
        Q = BQ[1];
        N = (A+Q)/2;
        D = foot(A,B,C);
        dot("$A$",A, A);
        dot("$B$",B, SW);
        dot("$C$",C, SW);
        dot("$M$",M, S);
        dot("$P$",P, NW);
        dot("$Q$",Q, SW);
        dot("$N$",N, NW);
        dot("$D$",D, SE);
        draw(A--B--C--A--M);
        draw(B--N--C--P);
        draw(c, dashed+red);
        draw(B--Q);
        draw(Q--A, blue);
        draw(A--D, blue);
        draw(M--N, blue);
        draw(rightanglemark(A,D,C,2), green);
        draw(rightanglemark(A,P,C,2), green);
        draw(circumcircle(triangle(C,P,D)), dashed);
    \end{asy}
\end{center}

Sekarang perhatikan dari \textit{power of a point} pada lingkaran $(ABPQ)$ dan $(CPDA)$ didapat
\begin{align*}
    BM \cdot MQ = PM \cdot MA = CM \cdot MD.
\end{align*}
Namun, karena $BM = CM$ maka didapatkan $MQ = MD$. Hal ini mengimplikasikan $QM : MD = QN : NA$ sehingga $\triangle MQN \sim \triangle DQA$ yang berakibat $\dangle QMN = \dangle QDA = 90^\circ$. Tentu saja karena $M$ titik tengah $BC$ dan $MN \perp BC$, maka $BN=CN$ seperti yang ingin dibuktikan.
\end{proof}

\begin{proof}[\textbf{Solusi 2.}] (Friday, 21 July 2023) Misalkan $A'$ adalah titik hasil pencerminan $A$ terhadap $M$. 
\begin{center}
    \begin{asy}
        import olympiad;
        import geometry;
        size(250);
        pair A,B,C,M,P,Q,N,A1;
        A=dir(75);
        B=dir(210);
        C=dir(-30);
        M=(B+C)/2;
        P=foot(C,A,M);
        A1 = M+(-1)*(A-M);
        circle c = circumcircle(triangle(A,B,P));
        circle d = circumcircle(triangle(A1,C,P));
        pair BQ[] = intersectionpoints(c, line(B,C));
        Q = BQ[1];
        N = (A+Q)/2;
        dot("$A$",A, NE);
        dot("$B$",B, SW);
        dot("$C$",C, E);
        dot("$M$",M, SE);
        dot("$P$",P, NW);
        dot("$Q$",Q, SW);
        dot("$N$",N, NW);
        dot("$A'$",A1, SE);
        draw(A--B--C--A--M--A1);
        draw(B--A1--C);
        draw(B--N--C--P);
        draw(c);
        draw(d, dashed+orange);
        draw(B--Q);
        draw(Q--A, blue);
        draw(A--A1, blue);
        draw(A1--Q, red);
        draw(N--M, red);
        draw(rightanglemark(C,Q,A1,2), green);
        draw(rightanglemark(A,P,C,2), green);
    \end{asy}
\end{center}

Perhatikan bahwa dengan \textit{power of a point} kita punya
\begin{align*}
    MQ \cdot MC = MQ \cdot MB = MP \cdot MA = MP \cdot MA'
\end{align*}
yang menunjukkan bahwa $Q,C,A',P$ konsiklis. Oleh karena itu $\dangle CQA' = \dangle CPA' = 90^\circ$ atau $QA' \perp BC$.

Di lain sisi, karena $AN=NQ$ dan $AM=MA'$, maka $NM \parallel QA'$ yang mengakibatkan $NM \perp BC$. Karena $M = \operatorname{BC}$, maka $BN=NC$. Terbukti.
\end{proof}